\chapter{Étude comparative}
\label{ch:etat-art}

\section{Périmètre de l'étude}

Cette étude positionne LexHub par rapport aux offres disponibles pour un cabinet marocain, selon la méthode retenue au premier suivi PFA : une \textbf{grille de critères} face aux applications existantes (cœur voix~$\rightarrow$~plaidoirie, gestion de cabinet, SaaS multi-tenant, équipe).

Les sources du benchmarking sont notamment : \texttt{smartavocat.ma}, \texttt{avocato.ma}, \texttt{avocatpro.ma}, \texttt{legaliq.ma}, \texttt{kanoun.com}, \texttt{adala.ai}, \texttt{diskod.ma}.

\section{Benchmarking --- offres existantes}
\label{sec:benchmark}

Le tableau~\ref{tab:benchmark} regroupe, pour chaque offre : les services principaux et ce qui manque face au besoin LexHub (flux plaidoirie vocale intégré au cabinet).

\begin{table}[H]
\centering
\caption{Benchmarking marché --- applications existantes vs besoin LexHub}
\label{tab:benchmark}
\footnotesize
\renewcommand{\arraystretch}{1.35}
\begin{tabular}{@{}p{2.4cm}p{4.6cm}p{5.6cm}@{}}
\toprule
\textbf{Offre} & \textbf{Services offerts} & \textbf{Ce qui manque face à LexHub} \\
\midrule
SmartAvocat \newline {\scriptsize smartavocat.ma} &
Gestion cabinet, dossiers bilingues &
Pas de flux plaidoirie vocale \\
\midrule
Avocato \newline {\scriptsize avocato.ma} &
Gestion + audiences / Mahakim &
Pas de parcours dictée $\rightarrow$ document métier \\
\midrule
AvocatPro \newline {\scriptsize avocatpro.ma} &
Gestion + Mahakim, organisation &
Pas de cœur plaidoirie \\
\midrule
Legaliq \newline {\scriptsize legaliq.ma} &
IA juridique / fiscale, documents &
Pas un SaaS cabinet complet \\
\midrule
Kanoun \newline {\scriptsize kanoun.com} &
IA droit marocain + OHADA &
Pas multi-tenant cabinet \\
\midrule
Adala \newline {\scriptsize adala.ai} &
Assistant IA, Q\&A, documents &
Pas d'équipe cabinet \\
\midrule
Dragon Legal / Diskod \newline {\scriptsize diskod.ma} &
Dictée desktop (Nuance) &
Pas SaaS, pas dossiers LexHub \\
\midrule
\textbf{LexHub} \newline {\scriptsize lexhub.space} &
Dictée $\rightarrow$ correction $\rightarrow$ document ; dossiers ; équipe ; SaaS multi-tenant &
--- \\
\bottomrule
\end{tabular}
\end{table}

\section{Positionnement de LexHub}

LexHub se situe à l'intersection :
\begin{itemize}[leftmargin=1.4em]
  \item d'un \textbf{SaaS cabinet} (inscription, isolation, équipe, dossiers, planning, activité) ;
  \item d'un \textbf{atelier de plaidoirie} (cinq étapes guidées : dictée, transcription/correction, réglages, document, export) ;
  \item d'un \textbf{export documentaire} (en-tête, typographie arabe, Word et PDF).
\end{itemize}

Le choix d'une \textbf{IA d'assistance} (transcription + structuration contrôlée) plutôt que d'une IA « auteur » est un parti pris scientifique et professionnel : on maximise la fidélité à la dictée, on refuse l'hallucination juridique. En une phrase : \emph{dicter $\rightarrow$ corriger $\rightarrow$ modifier $\rightarrow$ exporter (PDF/DOC) + dossiers + équipe}.
